\section{Test 1: conjunts numèrics, ordre i intervals}
\iniciTest{r1}{c,a,b,c,d,a,d,c,d,d}

\begin{enumerate}
\pregunta Quina de les següents igualtats és vertadera?

\begin{enumerate}
\opcio{a}{$\mathbb{N}\cup \mathbb{Q}=\mathbb{Z}$
}
\opcio{b}{$\mathbb{N}\cup \mathbb{Z}=\mathbb{Q}$
}
\opcio{c}{$\mathbb{I}\cap \mathbb{Q}=\varnothing $
}
\opcio{d}{$\mathbb{R}\cap \mathbb{Q}=\mathbb{I}$}
\end{enumerate}
\feedback


on $\mathbb{I}$ denota el conjunt dels nombres irracionals.

\pregunta Quina de les següents afirmacions és vertadera?

\begin{enumerate}
\opcio{a}{$%
\begin{array}{ccc}
x\in \mathbb{Q},y\notin \mathbb{Q} & \Longrightarrow & x+y\notin \mathbb{Q}%
\end{array}%
$
}
\opcio{b}{$%
\begin{array}{ccc}
x\notin \mathbb{Q},y\notin \mathbb{Q} & \Longrightarrow & x+y\notin \mathbb{Q%
}%
\end{array}%
$
}
\opcio{c}{$%
\begin{array}{ccc}
x\in \mathbb{Q},y\notin \mathbb{Q} & \Longrightarrow & x\cdot y\notin 
\mathbb{Q}%
\end{array}%
$
}
\opcio{d}{$%
\begin{array}{ccc}
x\notin \mathbb{Q},y\notin \mathbb{Q} & \Longrightarrow & x\cdot y\notin 
\mathbb{Q}%
\end{array}%
$}
\end{enumerate}
\feedback


\pregunta Quin dels següents nombres és irracional?

\begin{enumerate}
\opcio{a}{$3.1415$
}
\opcio{b}{$1.41421356...$
}
\opcio{c}{$2.1333...$
}
\opcio{d}{$0.777...$}
\end{enumerate}
\feedback


\pregunta Quantes aproximacions successives de nombres
racionals són necessàries per determinar un nombre irracional?

\begin{enumerate}
\opcio{a}{Basta amb dos, una per defecte i una altra per excés, suficientment
precises
}
\opcio{b}{Tantes com xifres decimals del nombre irracional es vulguin
conèixer
}
\opcio{c}{Infinites
}
\opcio{d}{Moltes però sempre en un nombre determinat}
\end{enumerate}
\feedback


\pregunta Quina de les següents afirmacions és falsa?

\begin{enumerate}
\opcio{a}{$%
\begin{array}{ccc}
a<b,c<d & \Longrightarrow & a+c<b+d%
\end{array}%
$
}
\opcio{b}{$%
\begin{array}{ccc}
a<b,c>d & \Longrightarrow & a-c<b-d%
\end{array}%
$
}
\opcio{c}{$%
\begin{array}{ccc}
a<b,c>0 & \Longrightarrow & ac<bc%
\end{array}%
$
}
\opcio{d}{$%
\begin{array}{ccc}
a<b,c<0 & \Longrightarrow & ac<bc%
\end{array}%
$}
\end{enumerate}
\feedback


\pregunta Si $A=\{\frac{(-1)^{n}}{n}:n\in \mathbb{N}\}$, aleshores
quina de les següents afirmacions és falsa?

\begin{enumerate}
\opcio{a}{$\sup A=1$
}
\opcio{b}{$\max A=\frac{1}{2}$
}
\opcio{c}{$\inf A=-1$
}
\opcio{d}{$\min A=-1$}
\end{enumerate}
\feedback


\pregunta Si $A=\{x\in \mathbb{R}:x^{2}-x<0\}$, aleshores

\begin{enumerate}
\opcio{a}{$A=(-\infty ,0]\cup \lbrack 1,+\infty )$
}
\opcio{b}{$A=[0,1]$
}
\opcio{c}{$A=\mathbb{R}-[0,1]$
}
\opcio{d}{$A=(0,1)$}
\end{enumerate}
\feedback


\pregunta Si $A=\{x\in \mathbb{R}:\left\vert x-1\right\vert <2\}$ i $B=\{x\in
\mathbb{R}:x^{2}-4x+3\geq 0\}$, aleshores

\begin{enumerate}
\opcio{a}{$A\cap B=(-\infty ,1]$
}
\opcio{b}{$A\cap B=[-1,1]$
}
\opcio{c}{$A\cup B=\mathbb{R}$
}
\opcio{d}{$A\cup B=[1,3]$}
\end{enumerate}
\feedback


\pregunta Quina de les següents igualtats és vertadera?

\begin{enumerate}
\opcio{a}{$(-4,2)=E_{1}(-1)$
}
\opcio{b}{$(-\infty ,3]\cap (0,2]=[0,2]$
}
\opcio{c}{$E_{2}(0)\cup (2,+\infty )=(-2,+\infty )$
}
\opcio{d}{$(3,15)=E_{6}(9)$}
\end{enumerate}
\feedback


\pregunta Quin dels següents fets és cert?

\begin{enumerate}
\opcio{a}{L’equació $(x^{2}-5)(16x^{2}-9)=0$ admet quatre solucions en $%
\mathbb{Q}$
}
\opcio{b}{$300$ és una aproximació per defecte de $367.19528$ amb 3 xifres
significatives
}
\opcio{c}{$1.99$ és una aproximació per arrodoniment fins a les centèsimes
de $1.98146$
}
\opcio{d}{Si $a$ i $b$ són nombres positius i $a^{2}<b^{2}$, aleshores $a<b$}
\end{enumerate}
\feedback

\end{enumerate}